\chapter{Options: Parity and Greeks}\label{ch:iv2:options-parity-and-greeks}

\begin{ivtrains}
What an option position pays, what no price may violate, what each Greek measures and where it is largest, and what a hedger's day looks like. Unless a stem says otherwise, options are European on a stock paying no dividends, one contract covers 100 shares, and Greeks are those of the Black-Scholes model.
\end{ivtrains}

\section{Payoffs and bounds}

\begin{question}[$\star$ \iqroles{QT, Trader, QR, Risk} \iqfirm{any} \iqlevel{junior}]\label{iq:iv2:options-parity-and-greeks:1}
You buy a call struck at 100 for \$5 and sell a call struck at 110 for \$2, on the same stock with the same expiry. What are your largest gain, your largest loss and your break-even at expiry? What view is this?
\end{question}

\begin{question}[$\star$ \iqroles{QT, Trader, QR, Risk} \iqfirm{any} \iqlevel{junior}]\label{iq:iv2:options-parity-and-greeks:2}
You buy a call and a put, both struck at 100, for \$8 in total. Where do you break even at expiry? What do you make or lose if the stock finishes at 105, and at 120? What are you betting on?
\end{question}

\begin{question}[$\star$ \iqroles{QT, Trader, QR} \iqfirm{market maker, hedge fund, bank} \iqlevel{junior}]\label{iq:iv2:options-parity-and-greeks:3}
You buy a 90 call, sell two 100 calls and buy a 110 call, paying \$2 in total. What is your largest profit, your largest loss, and where do you break even at expiry? Why can this position never have a negative price?
\end{question}

\begin{question}[$\star$ \iqroles{QT, Trader, QR} \iqfirm{market maker, hedge fund, bank} \iqlevel{junior}]\label{iq:iv2:options-parity-and-greeks:4}
Rates are zero. You buy a call and sell a put at the same strike of 100, for a net cost of \$2. What position does that make? What must the stock's price be, for there to be no arbitrage? What is your profit if the stock finishes at 110?
\end{question}

\begin{question}[$\star$ \iqroles{QT, Trader, QR} \iqfirm{market maker, hedge fund, bank} \iqlevel{junior}]\label{iq:iv2:options-parity-and-greeks:5}
A stock is at 100, the strike is 90, the time to expiry is one year and the rate is 5\%, compounded continuously. What is the lowest possible price of the European call? It trades at \$12. Give the trades that take the arbitrage, and what they make.
\end{question}

\begin{question}[$\star$ \iqroles{QT, Trader, QR} \iqfirm{market maker, hedge fund, bank} \iqlevel{junior}]\label{iq:iv2:options-parity-and-greeks:6}
The rate is 5\%, compounded continuously. A one-year European put struck at 100 trades at \$97. Why is it too dear, whatever the stock does? What is the highest it can be? What trades lock in the profit?
\end{question}

\begin{question}[$\star$ \iqroles{QT, Trader, QR} \iqfirm{market maker, hedge fund} \iqlevel{junior}]\label{iq:iv2:options-parity-and-greeks:7}
On one stock and one expiry, a put struck at 100 costs \$6, and a put struck at 110 costs \$5. Rates are zero. What is wrong? What do you trade, and what is your worst outcome? What is the largest gap the two puts could show?
\end{question}

\begin{question}[$\star$ \iqroles{QT, Trader, QR, Risk} \iqfirm{any} \iqlevel{junior}]\label{iq:iv2:options-parity-and-greeks:8}
For a long call and a long put, what are the signs of delta, gamma, vega and theta in the Black-Scholes model with zero rates? Which are the same for the two, and why must they be?
\end{question}

\begin{question}[$\star$ \iqroles{QT, Trader, Risk} \iqfirm{any} \iqlevel{junior}]\label{iq:iv2:options-parity-and-greeks:9}
You are long 10 call contracts with a delta of 0.6, short 5 put contracts with a delta of $-0.3$, and short 200 shares. Each contract covers 100 shares. What is your total delta, in shares? How many shares do you trade to make it zero?
\end{question}

\begin{question}[$\star$ \iqroles{QT, Trader} \iqfirm{market maker, hedge fund, bank} \iqlevel{junior}]\label{iq:iv2:options-parity-and-greeks:10}
You are long 5 at-the-money call contracts, each covering 100 shares, with a delta of 0.5 and a gamma of 0.02 per dollar. How many shares do you sell to hedge? The stock rises \$1. What is the new delta, and what do you do?
\end{question}

\begin{question}[$\star$ \iqroles{QT, Trader, Risk} \iqfirm{market maker, hedge fund, bank} \iqlevel{junior}]\label{iq:iv2:options-parity-and-greeks:11}
Your book has a gamma of 1\,000 shares per dollar and a theta of $-\$1\,500$ a day, and it is delta-hedged. Roughly what do you make or lose if the stock moves \$2 today? What about \$1? What daily move do you need to break even?
\end{question}

\begin{question}[$\star$ \iqroles{QT, Trader, Risk} \iqfirm{market maker, hedge fund, bank} \iqlevel{junior}]\label{iq:iv2:options-parity-and-greeks:12}
Your book has a vega of \$500 per point of volatility. The front month's implied volatility rises 3 points, and the far months' rises 1 point; the book's vega is \$300 in the front months and \$200 in the far ones. What is your profit? What would you make if you just used the whole vega and a single 3-point move?
\end{question}

\begin{question}[$\star$ \iqroles{QT, QR, Trader} \iqfirm{market maker, hedge fund, bank} \iqlevel{junior}]\label{iq:iv2:options-parity-and-greeks:13}
A one-year at-the-money call and put are on a \$100 stock with 20\% volatility and zero rates. By how much does each change, roughly, if the rate rises one percentage point? Which gains, and why?
\end{question}

\begin{question}[$\star$ \iqroles{QT, QR, Trader} \iqfirm{market maker, hedge fund, bank} \iqlevel{junior}]\label{iq:iv2:options-parity-and-greeks:14}
Why should you never exercise an American call early on a stock paying no dividends, when rates are positive? Prove it with a bound. Why does the same argument fail for a put?
\end{question}

\section{Parity under questioning}

\begin{question}[$\star\star$ \iqroles{QT, Trader, QR} \iqfirm{market maker, hedge fund, bank} \iqlevel{junior}]\label{iq:iv2:options-parity-and-greeks:15}
A stock is at 100, one-year options are struck at 100, and the rate is 5\%, compounded continuously, so the strike's present value is about 95.12. The call is at \$6 and the put at \$3. What should the difference be? Give the trades that make money, their cost today, and their result at each stock price at expiry.
\end{question}

\begin{question}[$\star\star$ \iqroles{QT, Trader, QR} \iqfirm{market maker, hedge fund} \iqlevel{junior, experienced}]\label{iq:iv2:options-parity-and-greeks:16}
A one-year call less a put at the strike 100 trades at \$3, on a stock at 100 that pays no dividend, with the strike's present value 95.12. Parity says \$4.88. You try to lock in the gap by buying the call, selling the put and shorting the stock, and your broker quotes a stock-borrow fee of 2\% a year. What does the fee do? What borrow fee does the gap imply? Why does this gap usually tell you about the borrow, not about an easy profit?
\end{question}

\begin{question}[$\star\star$ \iqroles{QT, QR} \iqfirm{market maker, hedge fund, bank} \iqlevel{junior, experienced}]\label{iq:iv2:options-parity-and-greeks:17}
For American options on a stock with no dividends, parity becomes a pair of inequalities. With stock at 100, strike at 100, a rate of 5\% compounded continuously and one year to expiry, between what limits must the call less the put lie? Prove the upper limit.
\end{question}

\begin{question}[$\star\star$ \iqroles{QT, Trader, QR} \iqfirm{market maker, hedge fund, bank} \iqlevel{junior, experienced}]\label{iq:iv2:options-parity-and-greeks:18}
You buy a call and sell a put at 90, and sell a call and buy a put at 110, all with one year left. What is that position worth at expiry, wherever the stock finishes? The market charges \$18.70 for it. What rate of interest, compounded continuously, does that imply? If you can borrow at 5\%, what do you do?
\end{question}

\begin{question}[$\star\star$ \iqroles{QT, QR} \iqfirm{market maker, hedge fund} \iqlevel{junior}]\label{iq:iv2:options-parity-and-greeks:19}
On a stock with no dividends and positive rates, a one-month European call struck at 100 is quoted at \$4, and a three-month one at \$3.50. Is that an arbitrage? What about with a large dividend in between? What is the trade if it is an arbitrage?
\end{question}

\begin{question}[$\star\star$ \iqroles{QT, QR, Risk} \iqfirm{market maker, hedge fund, bank} \iqlevel{junior, experienced}]\label{iq:iv2:options-parity-and-greeks:20}
A \$100 stock has 20\% volatility, and rates are zero. What is an at-the-money call's gamma with a year to go, and with one trading day to go, of 252 in a year? How much further do you have to go from the strike, in dollars, to halve the gamma with one day left? Why does a short-dated short option scare people?
\end{question}

\begin{question}[$\star\star$ \iqroles{QT, QR} \iqfirm{market maker, hedge fund, bank} \iqlevel{junior, experienced}]\label{iq:iv2:options-parity-and-greeks:21}
A European put struck at 100 has one year left on a \$50 stock with 20\% volatility and a 5\% rate, compounded continuously. What is it worth, and does it gain or lose value with one more day passing, all else equal? How much a year? When can theta be positive?
\end{question}

\begin{question}[$\star\star$ \iqroles{QT, Trader, QR} \iqfirm{market maker, hedge fund} \iqlevel{junior}]\label{iq:iv2:options-parity-and-greeks:22}
A stock is at 100, with 20\% volatility, zero rates and one year to expiry. What is the vega per volatility point for strikes of 80, 100 and 120? Why is it not symmetric around 100?
\end{question}

\begin{question}[$\star\star$ \iqroles{QT, QR, Trader} \iqfirm{market maker, hedge fund} \iqlevel{junior}]\label{iq:iv2:options-parity-and-greeks:23}
A \$100 stock has 40\% volatility, zero rates and a one-year option struck at 100. What are the option's delta, and the risk-neutral chance it finishes in the money? A trader says the delta is the probability. How far off is that here? When is it right?
\end{question}

\begin{question}[$\star\star$ \iqroles{QT, QR} \iqfirm{market maker, hedge fund, bank} \iqlevel{junior}]\label{iq:iv2:options-parity-and-greeks:24}
A one-year call on a stock with a dividend yield of 2\% has a delta of 0.55. What is the delta of the put at the same strike? Derive it. What is it with no dividend yield?
\end{question}

\begin{question}[$\star\star$ \iqroles{QT, Trader, QR} \iqfirm{market maker, hedge fund} \iqlevel{junior}]\label{iq:iv2:options-parity-and-greeks:25}
You own a straddle with a constant gamma of 0.04 per dollar, delta-hedged at 100. The stock goes to 103 and comes back to 100, and you rehedge at each of those points. What are your trades, and what is your gain per straddle? Check it with the formula for gamma. What must a day's theta be for this to break even?
\end{question}

\begin{question}[$\star\star$ \iqroles{QT, Trader, Risk} \iqfirm{market maker, hedge fund, bank} \iqlevel{junior, experienced}]\label{iq:iv2:options-parity-and-greeks:26}
This morning your book has a delta of 200 shares, a gamma of 50 shares per dollar, a vega of \$300 per point and a theta of $-\$150$ a day. Today the stock rises \$2 and implied volatility falls 1 point. What should you have made, line by line? You actually made \$120: what is unexplained, and what could explain it?
\end{question}

\begin{question}[$\star\star$ \iqroles{QT, Trader, QR} \iqfirm{market maker, hedge fund} \iqlevel{junior}]\label{iq:iv2:options-parity-and-greeks:27}
You sell a 90 put at \$3.00 and buy a 110 call at \$2.50. Their deltas are $-0.27$ and $0.35$, and their vegas are 0.31 and 0.37 per point. What do you make at expiry at 80, 100 and 120? What are your position's delta and vega? What do you gain if the 110 implied volatility rises one point and the 90 implied volatility falls one point?
\end{question}

\begin{question}[$\star\star$ \iqroles{QT, Trader, Analyst} \iqfirm{any} \iqlevel{junior}]\label{iq:iv2:options-parity-and-greeks:28}
Rates are zero. You own a stock at 100, buy a 90 put for \$3 and sell a 110 call for \$3. What are your largest gain and loss? Using put-call parity, show that the position is the same as a call spread plus some cash. What must the 90-110 call spread cost?
\end{question}

\begin{question}[$\star\star$ \iqroles{QT, Trader, Analyst} \iqfirm{any} \iqlevel{junior}]\label{iq:iv2:options-parity-and-greeks:29}
Rates are zero. You buy a stock at 100 and sell a 105 call for \$3. What is your profit at 85, 105 and 120? Which simpler position has the same payoff, and at what price must it trade? Why are people wrong to think selling calls against stock is ``income'' with no risk?
\end{question}

\begin{question}[$\star\star$ \iqroles{QT, QR, Trader} \iqfirm{bank, market maker, hedge fund} \iqlevel{junior}]\label{iq:iv2:options-parity-and-greeks:30}
A future trades at 100. European options on it are struck at 95, expire in half a year, and the rate is 5\%, compounded continuously. What is the call less the put? Derive it. Why does the formula look different from the one for options on a stock?
\end{question}

\begin{question}[$\star\star$ \iqroles{QT, QR, Trader} \iqfirm{market maker, hedge fund, bank} \iqlevel{junior, experienced}]\label{iq:iv2:options-parity-and-greeks:31}
An option pays \$1 if a stock finishes above 100, and nothing otherwise. The stock is at 100 with 20\% volatility, and rates are zero. What is it worth in the Black-Scholes model, and how do you get that from the price of ordinary calls? What is the hedger's problem near expiry?
\end{question}

\begin{question}[$\star\star$ \iqroles{QT, QR, Risk} \iqfirm{market maker, hedge fund} \iqlevel{junior}]\label{iq:iv2:options-parity-and-greeks:32}
You own a 100-110 call spread. Where is its gamma positive, where is it negative, and roughly where does it cross zero? Why is that? What does this mean for your hedging as the stock rises through the range?
\end{question}

\begin{question}[$\star\star$ \iqroles{QT, QR, Trader} \iqfirm{market maker, hedge fund} \iqlevel{junior, experienced}]\label{iq:iv2:options-parity-and-greeks:33}
A \$100 stock has one-year options struck at 100 and zero rates. The call's implied volatility is 25\%, and the put's is 27\%. In dollars, how far apart are the two options from parity? What is the implied borrow? Why do most systems show one implied volatility per strike?
\end{question}

\begin{question}[$\star\star$ \iqroles{QT, QR} \iqfirm{market maker, hedge fund, bank} \iqlevel{junior, experienced}]\label{iq:iv2:options-parity-and-greeks:34}
An American put struck at 100 is on a stock at \$1, with one year left and a 5\% rate, compounded continuously. Should you exercise now? Show it with a bound on what the European put can be worth. What does a positive rate have to do with it? Under a zero rate, would you ever exercise?
\end{question}

\section{Hedging and the harder cases}

\begin{question}[$\star\star\star$ \iqroles{QT, QR, Trader} \iqfirm{market maker, hedge fund, bank} \iqlevel{junior, experienced}]\label{iq:iv2:options-parity-and-greeks:35}
You have sold a \$1 million digital option: it pays \$1 million if a stock is above 100 at expiry, and nothing otherwise. The stock is at 100 with 20\% volatility, and rates are zero. What is your delta hedge, in shares, with a year to go, and with one trading day to go? A dealer usually prices such a thing as a call spread: how, and what does it do for the hedge?
\end{question}

\begin{question}[$\star\star\star$ \iqroles{QT, Trader, Risk} \iqfirm{market maker, hedge fund, bank} \iqlevel{junior, experienced}]\label{iq:iv2:options-parity-and-greeks:36}
It is the final day. You are short 1\,000 calls struck at 100, on 100 shares each, and long 100\,000 shares as a hedge. The stock closes at \$100.02, and holders may exercise after the close, in full or not at all, and you will not know until tomorrow. Over the evening the stock falls to \$98. What do you make or lose in each of the two cases? What is this risk called? What should you have done?
\end{question}

\begin{question}[$\star\star\star$ \iqroles{QT, QR, Trader} \iqfirm{market maker, hedge fund, bank} \iqlevel{junior, experienced}]\label{iq:iv2:options-parity-and-greeks:37}
An American call struck at 100 is on a stock at \$110 that pays \$2 tomorrow. After that no more dividends are due, and it expires in a quarter of a year. The rate is 5\%, compounded continuously. A European put at the same strike and expiry, after the dividend, is worth \$0.40. Should you exercise today? Derive a rule that compares the dividend with what you give up.
\end{question}

\begin{question}[$\star\star\star$ \iqroles{QT, QR} \iqfirm{market maker, hedge fund, bank} \iqlevel{junior, experienced}]\label{iq:iv2:options-parity-and-greeks:38}
A stock at 100 moves up 20\% or down 20\% in each of two periods. The one-period rate is 5\%, simple. Price a two-period put struck at 100, both European and American. Where do you exercise early, and what is the premium for that right?
\end{question}

\begin{question}[$\star\star\star$ \iqroles{QT, QR, Trader} \iqfirm{market maker, hedge fund} \iqlevel{junior, experienced}]\label{iq:iv2:options-parity-and-greeks:39}
You own a one-month call struck at 100, hedged daily, with an implied volatility of 30\% and zero rates. Its gamma is 0.046 per dollar when the stock is at 100, and 0.0038 when the stock is at 120. In one month, the stock moves 2\% a day in both cases, but in one it stays near 100, and in the other it jumps early to 120 and stays near there. Roughly, what does a typical day's gamma gain come to in each case? What does this say about realised volatility as a measure of your profit?
\end{question}

\begin{question}[$\star\star\star$ \iqroles{QT, QR, Trader} \iqfirm{market maker, hedge fund} \iqlevel{junior, experienced}]\label{iq:iv2:options-parity-and-greeks:40}
A stock is at 100, with 20\% volatility, zero rates and one year to expiry. You sell one at-the-money straddle and buy enough 90-110 strangles, a 90 put plus a 110 call, to have zero vega at 20\%. How many strangles? What do you collect? What is your profit if all implied volatilities move to 25\% or to 15\%, at once, with nothing else changing? What is the name of what you are long?
\end{question}

\begin{question}[$\star\star\star$ \iqroles{QR, QT} \iqfirm{market maker, hedge fund} \iqlevel{junior, experienced}]\label{iq:iv2:options-parity-and-greeks:41}
You sell an at-the-money one-year call on a \$100 stock with 20\% volatility and zero rates, and you hedge it 252 times, with volatility as implied. What is the standard deviation of your result? Give the approximation and check it by simulation. How much does that fall if you hedge four times as often? What does this mean for pricing?
\end{question}

\begin{question}[$\star\star\star$ \iqroles{QR, QT} \iqfirm{market maker, hedge fund} \iqlevel{experienced}]\label{iq:iv2:options-parity-and-greeks:42}
You hedge a short at-the-money one-year call on a \$100 stock at 20\% volatility, zero rates, every trading day of 252, paying 0.1\% of the value on each share traded, so 0.2\% round trip. By how much should you raise volatility to cover the costs, using Leland's rule? What does that cost in dollars of premium? Check it with a simulation of the rebalancing costs.
\end{question}

\begin{question}[$\star\star\star$ \iqroles{QT, QR, Trader} \iqfirm{market maker, hedge fund} \iqlevel{junior, experienced}]\label{iq:iv2:options-parity-and-greeks:43}
A \$100 stock has at-the-money options with a year and with a month to go, 20\% volatility on both, and zero rates. You are short one year-to-expiry lot. How many one-month lots do you buy to have zero vega? What is the net gamma then? If one-month volatility moves about $\sqrt{12}$ times as much as one-year volatility, how many one-month lots do you need instead?
\end{question}

\begin{question}[$\star\star\star$ \iqroles{QR, QT} \iqfirm{market maker, hedge fund} \iqlevel{junior, experienced}]\label{iq:iv2:options-parity-and-greeks:44}
A digital pays \$1 above 100 in a year, on a stock with 20\% volatility and zero rates. Does a rise in volatility make it worth more? Compute the change per point for stock prices of 90, 100 and 110, and find the price at which the sign changes. Why is it not at 100? What does an ordinary call have instead?
\end{question}

\begin{question}[$\star\star\star$ \iqroles{QT, Trader, Analyst} \iqfirm{market maker, hedge fund} \iqlevel{junior, experienced}]\label{iq:iv2:options-parity-and-greeks:45}
Dealers as a group are estimated to be long gamma of 1 million shares per dollar in a stock, near an expiry and near a large strike. The stock trades 5 million shares a day. If it rises \$1, what do the dealers do, and as a share of volume? What if they are short by the same amount? What happens to the price's behaviour in each case? What makes this estimate unreliable?
\end{question}

\begin{question}[$\star\star\star$ \iqroles{QR, QT} \iqfirm{market maker, hedge fund} \iqlevel{experienced}]\label{iq:iv2:options-parity-and-greeks:46}
Suppose a call's value $C(S, K)$ is homogeneous of degree one in the stock price and the strike. Show that $C = S\,\partial_S C + K\,\partial_K C$. Use this to get a delta from market prices alone, with zero rates, when calls at 99 and 101 with a \$100 stock are worth 8.44 and 7.52, and the call at 100 is worth 7.97. Check it against the Black-Scholes delta at 20\% volatility for a one-year call. What assumption about the smile does this rest on?
\end{question}

\begin{question}[$\star\star\star$ \iqroles{QT, QR} \iqfirm{market maker, hedge fund} \iqlevel{junior, experienced}]\label{iq:iv2:options-parity-and-greeks:47}
Rates are zero, and calls on one stock and expiry are quoted: 90 at \$14.00, 100 at \$8.00, 110 at \$5.50, and 120 at \$2.00. Which conditions do these prices meet, and which do they break? What is your trade, and your worst and best result at expiry? What does the table say about the chance of finishing above 95, above 105 and above 115?
\end{question}

\begin{question}[$\star\star\star\star$ \iqroles{QR} \iqfirm{hedge fund, market maker, bank} \iqlevel{experienced}]\label{iq:iv2:options-parity-and-greeks:48}
Show how to replicate any twice-differentiable payoff $f(S_T)$ with a bond, a forward and a strip of calls and puts, around a reference price $K_0$. Apply it to $f = S_T^2$ with a \$100 stock at 20\% volatility in the Black-Scholes model, zero rates and one year to expiry, with $K_0 = 100$. What does the strip come to, and what should it? What is the weight on each option?
\end{question}

\begin{question}[$\star\star\star\star$ \iqroles{QR, QT} \iqfirm{hedge fund, market maker} \iqlevel{experienced}]\label{iq:iv2:options-parity-and-greeks:49}
You sell a three-month at-the-money call on a \$100 stock at 30\% implied volatility, with zero rates, and the stock realises 20\%. What is your expected profit? Compare hedging with deltas at the implied volatility against hedging with deltas at the realised volatility. For each, what is the formula for the profit along a path, and what are the profit's mean and standard deviation with daily hedges? Which would you choose if you knew the 20\%, and why do most desks use the implied?
\end{question}

\begin{question}[$\star\star\star\star$ \iqroles{QR} \iqfirm{hedge fund, bank} \iqlevel{experienced}]\label{iq:iv2:options-parity-and-greeks:50}
A perpetual American put struck at 100 is on a stock with 20\% volatility and a 5\% rate, compounded continuously, and no dividends. Find its value and the stock price at which you exercise, by solving the equation it must satisfy. Show that the exercise rule is the one that maximises value. What is it worth with the stock at 100? What is the limit if the rate goes to zero?
\end{question}
